% ════════════════════════════════════════════
% 变量与估计方法（5.3.2）
% ════════════════════════════════════════════
\subsubsection{变量与估计方法}

数据整理阶段已把原始数据整理为 500 行、26 个参数的分析数据，本部分直接读取该文件，不再改动任何数值。因变量为 log1p 曝光量，设计矩阵的第一列为常数项，其余列依次为粉丝规模对数、帖子类型的 3 个哑变量、内容主题的 2 个哑变量、月份的 11 个哑变量、星期的 6 个哑变量、时段的 2 个哑变量与付费标记。全部哑变量采用 0/1 编码，参照类分别为 Photo、主题 1、1 月、周一与上午。

为保持系数解释口径，建模时不对连续变量做标准化，粉丝规模对数的系数仍表示粉丝规模每增加 1\% 对应的曝光变化。但在共线性诊断中，量纲差异会使条件数虚高，因此另构造标准化设计矩阵：把每一列减去其均值再除以其标准差，使所有列处于同一尺度。条件数只用于诊断，不参与回归估计，因而不会改变系数含义。

在确定设计矩阵后，需要选择与数据特征匹配的估计方法。曝光量重尾、存在异方差，普通最小二乘的经典标准误会低估参数不确定性；若直接使用稳健回归或广义线性模型，又会牺牲系数的直观解释。针对这一任务，需在经典最小二乘、HC3 稳健最小二乘、Huber 稳健回归、Gamma 广义线性模型与树模型之间择优，从估计目标、抗异方差能力与可解释性三个维度对比其适用性，如表\ref{tab:algo3}所示。

\begin{longtable}{>{\centering\arraybackslash}p{0.20\textwidth}>{\centering\arraybackslash}p{0.24\textwidth}>{\centering\arraybackslash}p{0.24\textwidth}>{\centering\arraybackslash}p{0.20\textwidth}}
  \caption{回归估计方法对比}
  \label{tab:algo3} \\
  \toprule
  方法 & 主要优点 & 主要局限 & 适用性 \\
  \midrule
  \endfirsthead
  \caption{回归估计方法对比（续）} \\
  \toprule
  方法 & 主要优点 & 主要局限 & 适用性 \\
  \midrule
  \endhead
  \bottomrule
  \endfoot
  经典最小二乘 & 系数含义直观，计算稳定 & 异方差下标准误有偏 & 作为基准与系数点估计 \\
  HC3 稳健最小二乘 & 不改变点估计，修正标准误 & 小样本下略偏保守 & 主推断方法 \\
  Huber 稳健回归 & 降低极端残差影响 & 损失函数改变，解释稍复杂 & 稳健性对照 \\
  Gamma 广义线性模型 & 匹配右偏计数，方差随均值增长 & 需假定条件分布 & 稳健性对照 \\
  梯度提升树 & 自动捕捉非线性与交互 & 系数不可解释，易过拟合 & 不适用于解释型任务 \\
\end{longtable}

如表\ref{tab:algo3}所示，区间估计的可靠性取决于标准误而非点估计，HC3 在不改变最小二乘点估计的前提下修正异方差导致的标准误偏差，最适合作为主推断方法\cite{white1980}。Huber 回归与 Gamma 模型从不同角度处理重尾，可作为结论是否依赖估计方法的对照\cite{huber1964,mccullagh1989}。树模型虽然预测灵活，但无法回答“付费带来多少增量”这类需要系数的问题，故不采用。综合以上分析，本部分选择 HC3 稳健最小二乘为主、Huber 与 Gamma 为辅的方案。
